\begin{center}
\textsc{Gauss} an \textsc{Olbers}.\quad Göttingen, 2.~April 1826.
\end{center}

\par\noindent\hspace*{\parindent}\mbox{.\ .\ .\ .\ .\ .}\quad Rücksichtlich meiner Messungen kann ich \textls[120]{für mich allein} wenig beschliessen.
Mein \textls[120]{Auftrag} ist im Grunde rücksichtlich des trigonometrischen Theils vollendet, und ich dachte, insofern \textsc{Schumacher} mich \textls[120]{gehörig} unterstützen will, im Spätsommer die Zenithsector-Beobachtungen vorzunehmen.
Der wankende Zustand meiner Gesundheit schreckt mich ab, auf erweiterte trigonometrische Messungen anzutragen; inzwischen habe ich vor kurzem, \textls[120]{unter uns gesagt}, in einem Schreiben an \textsc{Münster} erklärt, dass ich bereit bin, meine Kräfte auch noch künftig darauf zu wenden, falls solche gefordert werden sollten.

\par\noindent\hspace*{\parindent}Meine theoretischen Arbeiten lassen bei ihrem so \textls[120]{sehr} grossen Umfange leider noch viele Lücken; am leichtesten wäre mir geholfen, wenn ich mir
\par\clearpage
\noindent\null\hfill{\small\textbf{ZUR HANNOVERSCHEN TRIANGULATION.}}\hfill\makebox[0pt][r]{\fontsize{13}{16}\selectfont \textbf{377}}\par\medskip
\noindent erlaubte, mit der Bekanntmachung meiner Messungen zwar alle meine Rechnungseinrichtungen zu verbinden, aber deren \textls[120]{Ableitung} aus ihren höhern Gründen für ein ganz getrenntes Werk für glücklichere zukünftige Zeiten aufsparte.
Dann wäre nirgends ein Anstoss.
Vors erste werde ich die scharfe Ausgleichung meiner 32 Punkte, die 51 Dreiecke und 146 Richtungen liefern, vornehmen.
Die Höhenausgleichung (ein sehr viel leichteres Geschäft) habe ich in diesen Tagen vollendet.\quad\mbox{.\ .\ .\ .\ .\ .}
